\section{Characteristic surface and coincidence locus}

In the coordinates used by the application,
\[
U_- = \bar U-\frac{Y}{2}r(z),
\qquad
U_+ = \bar U+\frac{Y}{2}r(z),
\]
the zero-jump set is obtained by setting $Y=0$. We denote this
characteristic surface by
\[
\boxed{\mathcal C=\{(\tau,Y,z)\in\mathcal W:Y=0\}}.
\]
Thus, on $\mathcal C$,
\[
U_-=U_+=\bar U(\tau,z).
\]

Let
\[
DF(\bar U)=
\begin{pmatrix}
A&B\\
C&D
\end{pmatrix}.
\]
The discriminant of the characteristic polynomial is
\[
\Delta(\bar U)=(A-D)^2+4BC.
\]
The coincidence locus is the subset of the characteristic surface on which
the two characteristic speeds coincide. We denote it by
\[
\boxed{\mathcal E=\{(\tau,0,z)\in\mathcal C:\Delta(\bar U(\tau,z))=0\}}.
\]
Consequently,
\[
\boxed{\mathcal E\subset\mathcal C}.
\]

For fixed $z$, the implementation substitutes the fiber
\[
\bar U(\tau,z)=U_0(z)+\tau(-Q(z),P(z))^T
\]
into $\Delta$. Since $F$ is quadratic, $DF$ is affine in $\bar U$ and
$\Delta(\bar U(\tau,z))$ is a polynomial of degree at most two in $\tau$.
The function \texttt{coincidenceTauRoots} solves precisely this scalar
quadratic equation, while \texttt{coincidencePoints} samples the resulting
branches in the projective $z$ direction.

Away from $\mathcal E$, the characteristic surface has two spectral sheets.
If $s$ is the eigenvalue selected by $r(z)$ and
\[
\lambda_{\mathrm{other}}=\operatorname{tr}DF(\bar U)-s,
\]
we use
\[
\mathcal C_s=\{s<\lambda_{\mathrm{other}}\},
\qquad
\mathcal C_f=\{s>\lambda_{\mathrm{other}}\}.
\]
Their common spectral boundary is the coincidence locus $\mathcal E$, where
the two characteristic speeds agree.  This ordering is spectral and is not
determined by the sign of the projective coordinate $z$.

The notation is fixed throughout this project as follows:
\[
\mathcal C=\text{characteristic surface},
\qquad
\mathcal E=\text{coincidence locus}.
\]
